\scanpage{240}
- \mspage{185} - \margin{IX}
phisme (i.e. tel que $\underline O_S\longrightarrow f_*(\underline O_{S'})$ soit injectif) est un \uline{morphisme de descente effective} pour la catégorie fibrée des schémas abéliens relatifs.

b) Il semble que la méthode de démonstration de 13.5.2 doive permettre aussi de prouver le résultat analogue pour tout espace analytique $B$ propre sur $S$. Plus délicate sans doute est la question de savoir si, pour un espace analytique $X$ propre sur $S$, tel que le complété formel $\widehat X$ le long de $X_s$ soit algébrisable, $X$ est algébrisable sur $S$ au voisinage de $s$. Signalons seulement qu'il y a lieu ici de prendre la notion d'algébrisabilité relative de $X$ sur $S$ dans un sens plus large que dans [10], en entendant par là qu'au voisinage de $s$, $X$ provient d'un \uline{espace algébrique} sur $\underline O_{S,s}$ au sens de M. ARTIN [1].

13.5.4. Plaçons-nous dans les conditions de 13.5.1. Choisissons par ailleurs, pour tout entiers $n>0$, l'isomorphisme standart de groupes sur $\mathbb C$
\[
\mu_{n\mathbb C}\simeq(\mathbb Z/n\mathbb Z)_{\mathbb C}\ ,\quad m\longmapsto\exp(2i\pi m/n)\ ,
\]
ce qui nous donne, par passage à la limite sur les puissances d'un nombre fixé $\ell$, un isomorphisme canonique de faisceaux $\ell$-adiques sur $\mathbb C$ :
\[
\mathbb Z_\ell(1)_{\mathbb C}\simeq\mathbb Z_{\ell\mathbb C}\ ,
\]
que nous utiliserons pour identifier les deux membres.

De la suite exacte (13.1.3), on déduit, en appliquant la suite exacte des $\underline{\mathrm{Ext}}_{\mathbb Z}^i(\mathbb Z/n\mathbb Z,-)$, un isomorphisme canonique
\margin{\unsure{503|502}}
\margin{239}
\margin{503}
